\chapter{Bentuk Kuadratik}\label{ch:b2:quadratic}

Sebuah \omterm{def:b2:quadratic:def}{bentuk kuadratik} adalah bayangan aljabar sebuah geometri: bagian
bersignatur nol memipih, bagian positif melengkung ke satu arah, yang
negatif ke arah lain. Bab ini mereduksi setiap \omterm{def:b2:quadratic:def}{bentuk kuadratik} real
menjadi jumlah kuadrat bertanda $\pm$ (Gauss), dan membuktikan bahwa cacah
tandanya bersifat hakiki (Sylvester), lalu memahkotai geometri Euklides
dengan \emph{teorema spektral}: \omterm{def:b2:quadratic:adjoint}{endomorfisma simetrik} \omterm{def:b2:reduction:diag}{dapat
didiagonalkan} dalam basis ortonormal --- yakni teorema aljabar linear
terapan yang paling sering dipakai.

\section{Bentuk bilinear dan bentuk kuadratik}

\begin{definition}\label{def:b2:quadratic:def}
Sebuah \emph{bentuk bilinear \omterm{def:b2:quadratic:adjoint}{simetrik}} pada ruang vektor real $E$ adalah
pemetaan bilinear $\varphi \colon E \times E \to \R$ dengan $\varphi(x, y) =
\varphi(y, x)$; sedangkan \emph{bentuk kuadratik}\index{bentuk kuadratik}
yang bersesuaian adalah $q(x) = \varphi(x, x)$. Bentuk
$\varphi$ dipulihkan dari $q$ lewat
\emph{polarisasi}\index{kesamaan polarisasi}:
\[
\varphi(x, y) = \tfrac12\bigl(q(x + y) - q(x) - q(y)\bigr).
\]
Dalam basis $(e_i)$, \emph{matriks} bagi $\varphi$ adalah matriks
\omterm{def:b2:quadratic:adjoint}{simetrik} $B = (\varphi(e_i, e_j))$, dengan $q(x) = X^{\mathsf T} B X$;
lalu penggantian basis bermatriks $P$ mengganti $B$ dengan $P^{\mathsf T} B P$
(inilah \emph{kekongruenan}\index{matriks kongruen} --- bukan
keserupaan!). Adapun \emph{rank} bagi $q$ adalah $\operatorname{rk} B$
(yang invarian, sebab kekongruenan mengalikan dengan matriks terbalikkan).
\end{definition}

\begin{example}
Pada $\R^2$: $q(x, y) = x^2 + 4xy + y^2$ bermatriks $\begin{pmatrix}
1 & 2\\ 2 & 1\end{pmatrix}$. Sebuah hasil kali dalam persis merupakan
bentuk bilinear \omterm{def:b2:quadratic:adjoint}{simetrik} yang \omterm{def:b2:quadratic:def}{bentuk kuadratiknya} definit positif;
sedangkan bab ini mengkaji kasus umum yang tandanya tak tentu.
\end{example}

\begin{example}[Kekongruenan beraksi]\label{ex:b2:quadratic:congruence}
Ambillah $q(x, y) = x^2 + 4xy + y^2$ (bermatriks $B = \begin{pmatrix} 1
& 2\\ 2 & 1\end{pmatrix}$) beserta basis barunya $e_1' = (1, 1)$,
$e_2' = (1, -1)$, yakni $P = \begin{pmatrix} 1 & 1\\ 1 &
-1\end{pmatrix}$. Maka
\[
P^{\mathsf T}BP
= \begin{pmatrix} 1 & 1\\ 1 & -1\end{pmatrix}
\begin{pmatrix} 1 & 2\\ 2 & 1\end{pmatrix}
\begin{pmatrix} 1 & 1\\ 1 & -1\end{pmatrix}
= \begin{pmatrix} 6 & 0\\ 0 & -2\end{pmatrix} :
\]
dalam koordinat $(u, v)$ sepanjang basis barunya, $q = 6u^2 -
2v^2$ --- periksa: $x = u + v$, $y = u - v$ memberi $x^2 + 4xy +
y^2 = 6u^2 - 2v^2$ secara langsung. Perhatikan bahwa entri diagonal
barunya $6, -2$ \emph{bukan} \omterm{def:b2:reduction:eigen}{nilai eigen} $3, -1$ milik $B$: sebab
\omterm{def:b2:quadratic:def}{kekongruenan} menskalakan ulang, dan hanya keserupaan yang memelihara
\omterm{def:b2:reduction:eigen}{spektrum} --- tetapi tandanya sepakat, sebagaimana dituntut teorema
Sylvester. (Basis di sini ortogonal tetapi tak ortonormal; dan
menormalkannya dengan $\frac{1}{\sqrt2}$ akan membagi diagonalnya
dengan $2$ lalu memulihkan \omterm{def:b2:reduction:eigen}{nilai eigennya}.)
\end{example}

\begin{example}[Determinan Gram mengukur luas]\label{ex:b2:quadratic:gramarea}
Untuk $v_1, v_2$ pada ruang Euklides, matriks Gram $G =
\bigl(\langle v_i, v_j\rangle\bigr)$ mengemas panjang dan
sudutnya; sedangkan \omterm{def:b2:linalg:det}{determinannya} mengemas \emph{luas}:
\[
\det G = \norm{v_1}^2\norm{v_2}^2 - \langle v_1, v_2\rangle^2
= \norm{v_1}^2\norm{v_2}^2\bigl(1 - \cos^2\theta\bigr)
= \bigl(\norm{v_1}\,\norm{v_2}\sin\theta\bigr)^2 ,
\]
yakni kuadrat luas jajaran genjang yang bertumpu pada $v_1, v_2$ ---
dan Cauchy--Schwarz persis merupakan pernyataan $\det G \geq 0$.
Contoh yang dikerjakan: $v_1 = (1, 2, 2)$, $v_2 = (2, 1, -2)$ dalam
$\R^3$:
\[
G = \begin{pmatrix} 9 & 0\\ 0 & 9 \end{pmatrix},
\qquad
\det G = 81 :
\]
jadi vektornya ortogonal dan berpanjang $3$, sehingga merentang
jajaran genjang (di sini sebuah bujur sangkar) berluas $\sqrt{81} = 9$.
Pelajaran penutupnya: tak ada hasil kali silang dan tak ada sihir
berdimensi $3$ yang dipakai --- sebab $\sqrt{\det G}$ mengukur volume
berdimensi $k$ pada \emph{sembarang} dimensi, dan itulah titik tolak
Bagian I soal akhir pekan sekaligus integral luas permukaan yang
menyusul pada jilid ini.
\end{example}

\section{Reduksi Gauss dan inersia Sylvester}

\begin{theorem}[Reduksi Gauss]\label{thm:b2:quadratic:gauss}
Setiap \omterm{def:b2:quadratic:def}{bentuk kuadratik} $q$ pada ruang real berdimensi hingga dapat
dituliskan\index{reduksi Gauss}
\[
q = \sum_{i=1}^{s} \ell_i^2 - \sum_{j=1}^{t} m_j^2 ,
\]
dengan $\ell_1, \dots, \ell_s, m_1, \dots, m_t$ berupa bentuk linear
yang bebas linear; setara dengan itu, suatu basis membuat
matriks $q$ menjadi diagonal berentri $+1$ (sebanyak $s$ kali), $-1$
(sebanyak $t$ kali), dan $0$.
\end{theorem}

\begin{proof}
Secara induktif atas banyaknya peubah, dalam koordinat: $q(x_1,
\dots, x_n)$.

\emph{Kasus 1: suatu kuadrat muncul}, katakanlah koefisien $a$ bagi
$x_1^2$ tak nol. Kumpulkan semua suku ber-$x_1$ lalu lengkapkan
kuadratnya:
\[
q = a\Bigl(x_1 + \frac{1}{a}\,\lambda(x_2, \dots,
x_n)\Bigr)^{\!2} + q_1(x_2, \dots, x_n),
\]
dengan $\lambda$ linear dan $q_1$ kuadratik dalam peubah sisanya:
jadi satu bentuk bebas terpisah (ia melibatkan $x_1$, sedangkan
yang lain tidak), lalu induksinya berlaku pada $q_1$, dan tanda $\pm$-nya
datang dari tanda $a$ setelah diskalakan ulang dengan $\sqrt{\abs a}$.

\emph{Kasus 2: tak ada kuadrat, tetapi ada suku silang}, katakanlah
$b\,x_1x_2$ dengan $b \neq 0$. Pakailah kesamaan
\[
x_1x_2 = \tfrac14\bigl((x_1 + x_2)^2 - (x_1 - x_2)^2\bigr)
\]
setelah pengelompokan: dengan menulis $q = b\,x_1x_2 + x_1\alpha +
x_2\beta + q_2$ (dengan $\alpha, \beta, q_2$ dalam peubah lainnya), kita periksa
\[
q = \frac{b}{4}\Bigl[\Bigl(x_1 + x_2 + \frac{\alpha +
\beta}{b}\Bigr)^{2} - \Bigl(x_1 - x_2 + \frac{\beta -
\alpha}{b}\Bigr)^{2}\Bigr] + \widetilde q ,
\]
dengan $\widetilde q$ bebas dari $x_1, x_2$: jadi dua bentuk bebas
terpisah, dan induksinya merampungkannya.

Kebebasan bentuk yang terkumpul: urutkanlah kelompoknya sebagaimana
dihasilkan. Bentuk kelompok pertamanya memuat $x_1$ (Kasus 1)
atau $x_1, x_2$ (Kasus 2); sedangkan semua bentuk berikutnya bebas dari
peubah itu. Andaikan sebuah kombinasi linear semua bentuk terkumpulnya
lenyap. Membaca koefisien $x_1$ (dan $x_2$): hanya
kelompok pertamanya yang menyumbang, dan di dalam kelompok itu satu atau
dua bentuknya kasatmata bebas ($\ell$ sendirian; atau $\ell \pm
m$ dengan $\ell, m$ bebas): jadi koefisien kelompok pertamanya
lenyap. Kupas kelompok itu lalu ulangi: secara induktif sepanjang
kelompoknya, semua koefisiennya lenyap --- jadi seluruh keluarganya bebas,
yakni ketrianguleran yang dibuat gamblang.
\end{proof}

\begin{theorem}[Hukum inersia Sylvester]\label{thm:b2:quadratic:sylvester}
Pasangan $(s, t)$ pada \cref{thm:b2:quadratic:gauss} hanya bergantung
pada $q$, bukan pada reduksinya: ia disebut
\emph{signatur}\index{signatur bentuk kuadratik} bagi $q$.
Lebih jauh
\[
s = \max\{\dim F : q|_F \text{ definit positif}\},
\]
dan secara simetris untuk $t$.
\end{theorem}

\begin{proof}
Misalkan $q = \sum_{i \leq s}\ell_i^2 - \sum_{j\leq t} m_j^2$ lalu misalkan
$F_+$ rentang vektor pradual yang di atasnya $(\ell_i)$
membatas menjadi koordinat --- secara konkret: lengkapilah keluarga
bebas $(\ell_1, \dots, \ell_s, m_1, \dots, m_t)$
menjadi \omterm{def:b2:linalg:dual}{basis dual} $E^*$, lalu misalkan $(u_1, \dots, u_n)$
basis $E$ yang bentuk koordinatnya adalah itu (yakni basis
pradualnya: $\ell_i(u_k) = \delta_{ik}$ untuk $k \leq s$,
sedangkan bentuk berikutnya lenyap pada vektor sebelumnya). Tetapkan $F_+ =
\operatorname{Vect}(u_1, \dots, u_s)$: maka untuk $x = \sum_{i\leq
s}x_iu_i \in F_+$,
\[
\ell_i(x) = x_i,
\qquad m_j(x) = 0,
\qquad\text{jadi}\qquad
q(x) = \sum_{i\leq s}x_i^2 > 0 \quad (x \neq 0) :
\]
jadi $q|_{F_+}$ definit positif dan maksimum pada displainya
$\geq s$. Sebaliknya, misalkan $F$ sembarang subruang dengan $q|_F$
definit positif, dan $G = \{x : \ell_1(x) = \dots = \ell_s(x) =
0\}$, yang berkodimensi $\leq s$; pada $G$ berlaku $q(x) = -\sum m_j^2 \leq 0$.
Maka $F \cap G = \{0\}$ (sebab vektor tak nol di sana akan punya $q > 0$
sekaligus $q \leq 0$), jadi $\dim F \leq \dim E - \dim G \leq s$. Karena itu
maksimumnya sama dengan $s$ bagi \emph{setiap} reduksi: jadi $s$ bersifat
hakiki, dan $t = \operatorname{rk} q - s$ demikian pula.
\end{proof}

\begin{example}\label{ex:b2:quadratic:gaussexample}
$q(x, y, z) = xy + yz + zx$ (tanpa kuadrat). Dengan $x_1 = x$ dan $x_2 =
y$: $q = xy + z(x + y)$, lalu kesamaan dua kuadratnya memberi
\[
q = \tfrac14(x + y + 2z)^2 - \tfrac14(x - y)^2 - z^2 ,
\]
(uraikan untuk memeriksanya). Tiga bentuk bebas: jadi signaturnya $(1, 2)$
dan ranknya $3$. Satu arah positif, dua arah negatif: yakni geometri
``kerucut cahaya'' milik bentuk ini.
\end{example}

\begin{example}[Sebuah bentuk merosot, direduksi selengkapnya]\label{ex:b2:quadratic:degenerate}
$q(x, y, z) = xy + yz$ pada $\R^3$: tanpa kuadrat, jadi Kasus 2 dengan
pengelompokan $q = y(x + z)$. Kesamaan dua kuadratnya pada hasil kali
bentuk bebas $y$ dan $x + z$ memberi
\[
q = \frac14\bigl(y + x + z\bigr)^2 -
\frac14\bigl(y - x - z\bigr)^2 .
\]
Kedua bentuk linear $y + x + z$ dan $y - x - z$ bersifat
bebas (sebab selisihnya $2(x+z)$ dan jumlahnya $2y$), jadi
Sylvester langsung membaca: signaturnya $(1, 1)$ dan ranknya $2$ ---
sehingga \emph{merosot}. Kernel bentuk polarnya ditemukan dengan
menyelesaikan $\varphi(v, \cdot) = 0$: dengan matriks
$\frac12\begin{pmatrix} 0&1&0\\ 1&0&1\\ 0&1&0\end{pmatrix}$,
kernelnya adalah $\{y = 0,\ x + z = 0\} = \R\,(1, 0, -1)$, yakni
arah yang membuat $q$ tak melihat apa pun. Pelajaran penutupnya:
kekurangan rank muncul pada Gauss sebagai ``kehabisan peubah''
--- sebab reduksinya hanya menghasilkan dua kuadrat dari tiga
dimensi, dan dimensi yang hilang itu persis kernelnya.
\end{example}

\begin{example}[Satu bentuk, dua jalan menuju signaturnya]\label{ex:b2:quadratic:tworoads}
$q(x, y, z) = 2x^2 + 2y^2 + 2z^2 + 2xy + 2yz$, bermatriks
$\begin{pmatrix} 2 & 1 & 0\\ 1 & 2 & 1\\ 0 & 1 & 2
\end{pmatrix}$. \emph{Jalan 1, Gauss:} lengkapkan kuadratnya secara
berurutan,
\[
q = 2\Bigl(x + \frac y2\Bigr)^{\!2} + \frac32 y^2 + 2yz + 2z^2
= 2\Bigl(x + \frac y2\Bigr)^{\!2}
+ \frac32\Bigl(y + \frac{2z}{3}\Bigr)^{\!2} + \frac43 z^2 :
\]
jadi tiga kuadrat positif pada bentuk yang bebas, sehingga signaturnya
$(3, 0)$: yakni definit positif. \emph{Jalan 2, \omterm{def:b2:reduction:eigen}{nilai eigen}:} matriksnya
adalah $2I + N$ yang tridiagonal dengan $N$ matriks tetangganya; dan
\omterm{def:b2:reduction:eigen}{nilai eigennya} $2 + \sqrt2$, $2$, $2 - \sqrt2$ (periksalah
\omterm{def:b2:reduction:eigen}{vektor eigen} $(1, \pm\sqrt2, 1)$ dan $(1, 0, -1)$), yang semuanya
positif: jadi vonisnya sama, menurut
\cref{cor:b2:quadratic:principalaxes}. Gauss lebih cepat;
sedangkan \omterm{def:b2:reduction:eigen}{nilai eigen} berkata lebih banyak (ia memberi sumbu utamanya dan
nilai ekstrem $q$ pada bolanya). Pelajaran penutupnya: pivot
positif $2, \frac32, \frac43$ milik Gauss persis merupakan
nisbah $\frac{\Delta_k}{\Delta_{k-1}}$ minor utama terdepannya
($\Delta_1 = 2$, $\Delta_2 = 3$, $\Delta_3 = 4$) --- dan soal
akhir pekan membuktikannya secara umum.
\end{example}

\begin{method}[Menghitung sebuah signatur: tiga rute]\label{met:b2:quadratic:signature}
\begin{enumerate}
  \item \emph{Gauss} (selalu berhasil dan tercepat dengan tangan):
        lengkapkan kuadratnya secara berurutan, lalu Kasus 2 bila
        tak ada kuadrat yang tersedia; lalu cacahlah tandanya. Periksa
        bahwa bentuk linear terkumpulnya bebas --- sebab bentuk yang
        lebih sedikit daripada peubahnya berarti ada
        kernel (\cref{ex:b2:quadratic:degenerate}).
  \item \emph{Minor terdepan} (untuk uji kedefinitan): semua
        $\Delta_k > 0$ bila dan hanya bila definit positif
        (\cref{exo:b2:quadratic:8}); dan pivot
        $\Delta_k/\Delta_{k-1}$ bahkan memberi koefisien
        Gaussnya (soal akhir pekan). Rute ini gagal secara diam-diam
        bila suatu $\Delta_k = 0$: maka kembalilah ke rute 1.
  \item \emph{\omterm{def:b2:reduction:eigen}{Nilai eigen}} (paling informatif, paling mahal):
        yakni tanda \omterm{def:b2:reduction:eigen}{spektrumnya}
        (\cref{cor:b2:quadratic:principalaxes}); ia juga memberi
        sumbu utamanya beserta nilai ekstrem $q$ pada
        bola satuannya. Pilihlah rute ini bila struktur eigennya
        toh diperlukan.
\end{enumerate}
\end{method}

\section{Teorema spektral}

Kini misalkan $E$ ruang Euklides (dengan hasil kali dalam
$\langle\cdot,\cdot\rangle$, lihat jilid Tahun ke-1).

\begin{definition}[Adjoin; endomorfisma simetrik]\label{def:b2:quadratic:adjoint}
Untuk $u \in \mathcal{L}(E)$, \emph{adjoin}\index{adjoin}
$u^*$ adalah satu-satunya endomorfisma dengan
\[
\langle u(x), y\rangle = \langle x, u^*(y)\rangle
\qquad (x, y \in E);
\]
dan dalam basis ortonormal berlaku $\operatorname{Mat}(u^*) =
\operatorname{Mat}(u)^{\mathsf T}$. Adapun $u$ disebut
\emph{simetrik}\index{endomorfisma simetrik} (atau swa-adjoin) apabila
$u^* = u$ --- setara dengan matriksnya dalam basis ortonormal bersifat
simetrik.
\end{definition}

\begin{proof}[Keberadaan dan ketunggalan adjoinnya]
Untuk $y$ yang tetap, bentuk $x \mapsto \langle u(x), y\rangle$ bersifat
linear, sehingga (berkat dimensi hingganya) berbentuk $\langle x,
z_y\rangle$ bagi $z_y$ yang tunggal --- sedangkan pemetaan $y \mapsto z_y =:
u^*(y)$ bersifat linear berkat ketunggalan itu. Adapun pengenalan matriksnya:
bacalah $\langle u(e_i), e_j\rangle$ dari kedua arah.
\end{proof}

\begin{example}[Adjoinnya bergantung pada hasil kali dalamnya]\label{ex:b2:quadratic:weightedadjoint}
Pada $\R^2$ ambillah hasil kali dalam \emph{berbobot} $\langle x,
y\rangle_D = x_1y_1 + 2x_2y_2$ (bermatriks $D =
\operatorname{diag}(1,2)$) dan $u$ bermatriks $A =
\begin{pmatrix} 0 & 1\\ 0 & 0\end{pmatrix}$ dalam basis
kanoniknya. Dari $\langle u(x), y\rangle_D = (Ax)^{\mathsf T}Dy =
x^{\mathsf T}(A^{\mathsf T}D)y$ dan $\langle x, u^*(y)\rangle_D
= x^{\mathsf T}(DA^*)y$, matriks \omterm{def:b2:quadratic:adjoint}{adjoinnya} adalah
\[
A^* = D^{-1}A^{\mathsf T}D
= \begin{pmatrix} 1 & 0\\ 0 & \tfrac12\end{pmatrix}
\begin{pmatrix} 0 & 0\\ 1 & 0\end{pmatrix}
\begin{pmatrix} 1 & 0\\ 0 & 2\end{pmatrix}
= \begin{pmatrix} 0 & 0\\ \tfrac12 & 0\end{pmatrix}
\neq A^{\mathsf T} .
\]
Periksa kewarasannya di $x = (1,0)$, $y = (0,1)$:
\[
\langle u(x), y\rangle_D = \langle (0,0), y\rangle_D = 0,
\quad
\langle x, u^*(y)\rangle_D
= \langle(1,0), (0,\tfrac12)\rangle_D = 0 ;
\]
lalu di $x = (0,1)$, $y = (1,0)$:
\[
\langle u(x), y\rangle_D = \langle(1,0),(1,0)\rangle_D = 1,
\quad
\langle x, u^*(y)\rangle_D =
\langle(0,1),(0,\tfrac12)\rangle_D = 1 .
\] Pelajaran penutupnya: ``$\operatorname{Mat}(u^*)
= \operatorname{Mat}(u)^{\mathsf T}$'' adalah pernyataan tentang
basis \emph{ortonormal} belaka; sebab secara umum metrik $D$
ikut campur, persis seperti pada reduksi serentak
soal akhir pekan.
\end{example}

\begin{theorem}[Teorema spektral]\label{thm:b2:quadratic:spectral}
Misalkan $u$ \omterm{def:b2:quadratic:adjoint}{endomorfisma simetrik} ruang Euklides $E$. Maka
$E$ punya \emph{basis ortonormal beranggotakan \omterm{def:b2:reduction:eigen}{vektor eigen}} $u$; semua
\omterm{def:b2:reduction:eigen}{nilai eigennya} real, dan ruang eigen bagi \omterm{def:b2:reduction:eigen}{nilai eigen} yang berbeda saling
ortogonal. Bentuk matriksnya: setiap matriks \omterm{def:b2:quadratic:adjoint}{simetrik} real $A$ dituliskan
\[
A = P\,D\,P^{\mathsf T},
\qquad P \text{ ortogonal } (P^{\mathsf T} P = I),\ D
\text{ diagonal}.\index{teorema spektral}
\]
\end{theorem}

\begin{proof}
\emph{Sebuah \omterm{def:b2:reduction:eigen}{vektor eigen} ada.} Fungsi $x \mapsto \langle u(x),
x\rangle$ \omterm{def:b2:metric:continuity}{kontinu} pada bola satuan $S$ milik $E$, yang bersifat
\omterm{def:b2:metric:compact}{kompak} (berkat dimensi hingganya, \cref{thm:b2:nvs:finitedim}): jadi ia
mencapai maksimumnya $\lambda$ di suatu $a \in S$. Kita klaim $u(a) = \lambda
a$. Untuk sembarang $y \perp a$ dengan $\norm y = 1$ dan $t \in \R$,
vektor $x_t = \frac{a + ty}{\sqrt{1 + t^2}}$ terletak pada $S$
(sebab $\norm{a + ty}^2 = 1 + t^2$ menurut Pythagoras); lalu menguraikan
fungsi yang dimaksimumkan itu,
\[
g(t) = \langle u(x_t), x_t\rangle
= \frac{\langle u(a), a\rangle + 2t\langle u(a), y\rangle +
t^2\langle u(y), y\rangle}{1 + t^2}
\]
(kesimetrikan $u$ menyatukan kedua suku silangnya: $\langle u(a),
y\rangle = \langle a, u(y)\rangle = \langle u(y), a\rangle$).
Fungsi $g$ terdiferensialkan terhadap $t$ dengan maksimum di $t =
0$; jadi aturan hasil bagi di $0$ memberi
\[
g'(0) = \frac{2\langle u(a), y\rangle\cdot 1 - \langle u(a),
a\rangle\cdot 0}{1} = 2\langle u(a), y\rangle = 0 .
\]
Jadi $u(a)$ ortogonal terhadap seluruh hiperbidang $a^\perp$: sehingga $u(a)
\in (a^{\perp})^{\perp} = \R a$, yakni $u(a) = \mu a$; dan $\mu =
\langle u(a), a\rangle = \lambda$.

\emph{Induksinya.} Komplemen ortogonal $F = a^\perp$ bersifat
stabil terhadap $u$: sebab untuk $x \perp a$ berlaku $\langle u(x), a\rangle = \langle x,
u(a)\rangle = \lambda\langle x, a\rangle = 0$. Pembatasan
$u|_F$ \omterm{def:b2:quadratic:adjoint}{simetrik} terhadap hasil kali dalam terimbasnya; jadi secara induktif atas
dimensinya, $F$ punya basis eigen ortonormal; lalu sisipkan $a$ di depannya.

\emph{Pelengkapnya.} \omterm{def:b2:reduction:eigen}{Nilai eigennya} adalah bilangan real $\langle u(e),
e\rangle$ pada basis eigennya. Keortogonalan \omterm{def:b2:reduction:eigen}{ruang eigennya}: $u(x) =
\lambda x$ dan $u(y) = \mu y$ memberi $\lambda\langle x, y\rangle =
\langle u(x), y\rangle = \langle x, u(y)\rangle = \mu\langle x,
y\rangle$, jadi $\langle x, y\rangle = 0$ bila $\lambda \neq \mu$.
Bentuk matriksnya: kolom $P$ = basis eigen ortonormalnya.
\end{proof}

\begin{example}[Sebuah jalannya spektral yang lengkap]\label{ex:b2:quadratic:fullrun}
Diagonalkan secara ortogonal $A = \begin{pmatrix} 1 & 2\\ 2 &
1\end{pmatrix}$. \omterm{def:b2:reduction:charpoly}{Polinomial karakteristiknya} $(1 - \lambda)^2 -
4$: jadi \omterm{def:b2:reduction:eigen}{nilai eigennya} $3$ dan $-1$. \omterm{def:b2:reduction:eigen}{Vektor eigennya}: $(A - 3I)v = 0$
memberi $v_1 = \frac{1}{\sqrt2}(1,1)$; sedangkan $(A + I)v = 0$ memberi $v_2
= \frac{1}{\sqrt2}(1,-1)$ --- yang ortogonal, sebagaimana dijamin
\cref{thm:b2:quadratic:spectral} tanpa
perhitungan. Dengan $P = (v_1\ v_2)$ (yakni perputaran sebesar
$\frac\pi4$):
\[
P^{\mathsf T}AP = \begin{pmatrix} 3 & 0\\ 0 & -1
\end{pmatrix},
\qquad
x^2 + 4xy + y^2 = 3u^2 - v^2
\quad\text{pada kerangka terputar} .
\]
Jadi bentuk pada \cref{exo:b2:quadratic:1} adalah bentuk bertipe
hiperbola: signaturnya $(1,1)$, konsisten dengan reduksi Gaussnya
$(x + 2y)^2 - 3y^2$ --- kuadrat yang berbeda, signatur yang sama,
sebagaimana dituntut Sylvester. Pelajaran penutupnya: Gauss memberi jawabannya
lebih cepat, tetapi rute spektralnya juga melaporkan bahwa pada lingkaran
satuan $q$ berjangkau persis $\intcc{-1}{3}$, yang dicapai sepanjang
$v_2$ dan $v_1$: jadi kerja tambahannya membeli geometri.
\end{example}

\begin{corollary}[Sumbu utama; uji kepositifan]\label{cor:b2:quadratic:principalaxes}
\begin{enumerate}
  \item Setiap \omterm{def:b2:quadratic:def}{bentuk kuadratik} $q$ pada ruang Euklides \omterm{def:b2:reduction:diag}{dapat
        didiagonalkan} dalam suatu basis \emph{ortonormal}: $q(x) =
        \sum_i \lambda_i x_i^2$ dengan $\lambda_i$ \omterm{def:b2:reduction:eigen}{nilai eigen} matriks
        \omterm{def:b2:quadratic:adjoint}{simetrik} milik $q$; sedangkan signaturnya mencacah \omterm{def:b2:reduction:eigen}{nilai eigen}
        yang positif dan yang negatif.
  \item Sebuah matriks \omterm{def:b2:quadratic:adjoint}{simetrik} bersifat semidefinit positif (masing-masing
        definit) bila dan hanya bila semua \omterm{def:b2:reduction:eigen}{nilai eigennya} $\geq 0$
        (masing-masing $> 0$); dan nilai ekstrem hasil bagi Rayleighnya
        lalu menjadi
        \[
        \min_{\norm x = 1} \langle Ax, x\rangle = \lambda_{\min},
        \qquad
        \max_{\norm x = 1} \langle Ax, x\rangle = \lambda_{\max} .
        \]
\end{enumerate}
\end{corollary}

\begin{proof}
(1) Tulislah $q(x) = \langle A x, x\rangle$ dengan $A$ \omterm{def:b2:quadratic:adjoint}{simetrik} (yakni
matriks $q$ dalam basis ortonormal); lalu diagonalkan $A$ lewat
teorema spektralnya: untuk $x = \sum x_ie_i$ dalam basis eigen
ortonormalnya,
\[
q(x) = \Bigl\langle \sum_i \lambda_ix_ie_i,\ \sum_j
x_je_j\Bigr\rangle = \sum_i \lambda_i x_i^2
\]
(keortonormalannya membunuh suku silangnya). Tanda $\pm$ milik
$\lambda_i$ mencacah signaturnya menurut Sylvester: sebab menskalakan ulang tiap
koordinat dengan $\sqrt{\abs{\lambda_i}}$ memamerkan sebuah \omterm{thm:b2:quadratic:gauss}{reduksi
Gauss} dengan bentuk yang bebas.

(2) Dalam basis eigennya, $\langle Ax, x\rangle = \sum \lambda_i
x_i^2$, yang terjepit antara $\lambda_{\min}\norm x^2$ dan
$\lambda_{\max}\norm x^2$, dengan kesamaan pada \omterm{def:b2:reduction:eigen}{vektor eigen} yang
bersesuaian; jadi kepositifan semua \omterm{def:b2:reduction:eigen}{nilai eigennya} setara
dengan kepositifan bentuknya.
\end{proof}

\begin{example}\label{ex:b2:quadratic:spectralexample}
$A = \begin{pmatrix} 2 & 1\\ 1 & 2 \end{pmatrix}$: \omterm{def:b2:reduction:eigen}{nilai eigennya} $3$
(dengan \omterm{def:b2:reduction:eigen}{vektor eigen} $\frac{1}{\sqrt2}(1,1)$) dan $1$
(dengan $\frac{1}{\sqrt2}(1,-1)$). \omterm{def:b2:quadratic:def}{Bentuk kuadratik} $2x^2 + 2xy + 2y^2$
menjadi $3X^2 + Y^2$ dalam kerangka ortonormal yang terputar: yakni
sumbu utama sebuah elips, yang terhitung. \omterm{thm:b2:quadratic:gauss}{Reduksi Gauss} juga mencapai
bentuk diagonal, tetapi hanya teorema spektral yang mencapainya \emph{tanpa
memuaikan panjang}.
\end{example}

\begin{example}[Sebuah elips yang dikenali sepenuhnya]\label{ex:b2:quadratic:ellipserun}
Kurva apakah $5x^2 + 4xy + 2y^2 = 6$? Matriks
$\begin{pmatrix} 5 & 2\\ 2 & 2\end{pmatrix}$ berpolinomial
karakteristik $\lambda^2 - 7\lambda + 6 = (\lambda - 1)(\lambda -
6)$: jadi \omterm{def:b2:reduction:eigen}{nilai eigennya} $1$ dan $6$, keduanya positif --- yakni sebuah elips.
\omterm{def:b2:reduction:eigen}{Vektor eigen} ortonormalnya: untuk $\lambda = 1$, selesaikan
$\begin{pmatrix} 4 & 2\\ 2 & 1\end{pmatrix}v = 0$ sehingga $v_1 =
\frac{1}{\sqrt5}(1, -2)$; sedangkan untuk $\lambda = 6$: $v_2 =
\frac{1}{\sqrt5}(2, 1)$. Dalam koordinat terputar $(X, Y)$
sepanjang $(v_2, v_1)$ persamaannya menjadi
\[
6X^2 + Y^2 = 6,
\qquad\text{yakni}\qquad
X^2 + \frac{Y^2}{6} = 1 :
\]
jadi setengah sumbunya $1$ (sepanjang $v_2$) dan $\sqrt6$ (sepanjang $v_1$).
Pelajaran penutupnya: bentuk kasarnya gratis --- sebab $\det
= 6 > 0$ dan trace yang positif sudah mengumumkan sebuah elips sebelum
sebuah \omterm{def:b2:reduction:eigen}{vektor eigen} pun dihitung --- tetapi hanya teorema spektral yang
menyerahkan arah dan panjang sumbunya, yakni geometrinya yang
sesungguhnya.
\end{example}

\begin{example}[Ekstrem pada bolanya, dibaca dari spektrumnya]\label{ex:b2:quadratic:sphereextremes}
Berapakah nilai ekstrem $q(x,y,z) = 2xy + 2yz + 2zx$ pada
bola satuannya? Matriksnya (yakni luar diagonal serba-satu pada
\cref{exo:b2:quadratic:2}) bernilai eigen $2$ dan $-1$
(ganda), jadi menurut \cref{cor:b2:quadratic:principalaxes}~(2):
\[
\max_{\norm v = 1} q(v) = 2
\ \text{ di } v = \tfrac{1}{\sqrt3}(1,1,1),
\qquad
\min_{\norm v = 1} q(v) = -1
\ \text{ pada lingkaran } x + y + z = 0 .
\]
Tanpa kalkulus dan tanpa pengganda Lagrange: teorema spektral
menyelesaikan pengoptimalan berkendala ini secara langsung --- lalu
memamerkan pemaksimumnya. Pelajaran penutupnya: bandingkan dengan
metode pengganda pada bab kalkulus diferensial, yang menemukan
titik kritis yang sama dengan kerja lebih banyak; sebab untuk fungsi
tujuan \emph{kuadratik} pada bola, spektrumlah jalan rajanya (dan soal
akhir pekan bab Hermitian membangun seluruh
teori Courant--Fischer di atas pengamatan ini).
\end{example}

\begin{remark}[Jebakan yang sering muncul]
\emph{(i) \omterm{def:b2:quadratic:def}{Kekongruenan} bukan keserupaan:} penggantian basis bagi
sebuah bentuk bekerja lewat $P^{\mathsf T}BP$, bukan $P^{-1}BP$; jadi
\omterm{def:b2:reduction:eigen}{nilai eigennya} \emph{bukan} invarian sebuah \omterm{def:b2:quadratic:def}{bentuk kuadratik} ($I$ dan $4I$
kongruen lewat $P = 2I$) --- hanya tandanya yang invarian
(Sylvester). Berbicaralah tentang \omterm{def:b2:reduction:eigen}{nilai eigen} sebuah bentuk hanya setelah
sebuah hasil kali dalam ditetapkan. \emph{(ii) Entri positif tak
membuktikan apa pun:} $\begin{pmatrix} 1 & 2\\ 2 & 1\end{pmatrix}$ punya
entri yang semuanya positif namun bersignatur $(1,1)$ (sebab $\det = -3$);
sebaliknya matriks definit positif boleh punya entri luar diagonal yang
negatif (\cref{ex:b2:quadratic:tworoads} yang digeser:
$2I - N$ sama saja berhasil). Pakailah
\cref{met:b2:quadratic:signature}. \emph{(iii) Kuadrat yang
bergantungan:} menulis $q = \ell_1^2 - \ell_2^2$ tak berkata apa pun bila
$\ell_1, \ell_2$ sebanding --- sebab $x^2 + 2xy + y^2 =
(x+y)^2$ berrank $1$, bukan $2$; jadi periksalah selalu kebebasannya
sebelum membaca signaturnya. \emph{(iv) Ekstrem pada bola
tanpa kekompakan:} batas Rayleigh pada
\cref{cor:b2:quadratic:principalaxes} tercapai karena
bolanya \omterm{def:b2:metric:compact}{kompak}; sedangkan pada bola \omterm{def:b2:metric:topology}{terbuka} atau pada seluruh ruangnya, sebuah
bentuk yang tak tentu tak punya maksimum maupun minimum.
\end{remark}

\begin{remark}[Di mana ini dipakai]
Teorema spektral adalah hasil buku ini yang paling banyak diekspor:
statistika mendiagonalkan matriks kovariansi dengannya
(yakni analisis komponen utama), mekanika menyarikan moda normal
osilasi darinya (yakni reduksi serentak pada
soal akhir pekan), analisis numerik membangun dekomposisi Cholesky dan
dekomposisi nilai singular di atasnya (pada soal yang sama), dan bab
berikutnya mengangkutnya ke ruang Hermitian kompleks. Adapun jilid
Tahun ke-3 membuktikan awataranya yang berdimensi tak hingga untuk operator
\omterm{def:b2:quadratic:adjoint}{swa-adjoin} yang \omterm{def:b2:metric:compact}{kompak}, tempat hujah kekompakan pada bukti berdimensi
hingga di atas menjadi seluruh ceritanya.
\end{remark}

\begin{remark}[Pandangan ke depan di dalam jilid ini]
\omterm{def:b2:quadratic:def}{Bentuk kuadratik} menjalari sisa Buku 4 dalam tiga
samaran. Sebagai \emph{Hessian}: bab kalkulus diferensial
menggolongkan titik kritis lewat signatur bentuk
orde duanya, jadi keinvarianan Sylvesterlah yang membuat
``pelana'' menjadi kata yang terdefinisi baik. Sebagai \emph{energi}:
osilator pada bab persamaan diferensial mengusung
energi kuadratik $\frac12x'^2 + \frac12\omega^2x^2$, dan
reduksi serentak pada soal akhir pekan bab ini
persis merupakan penyarian moda normalnya. Sebagai \emph{geometri}:
konik pada bab ini tumbuh menjadi permukaan kuadrik pada
bab geometri, tempat bentuk fundamental kedua sebuah
permukaan --- yakni \omterm{def:b2:quadratic:def}{bentuk kuadratik} pada tiap bidang singgungnya --- punya
signatur yang memutuskan apakah permukaannya melengkung seperti mangkuk
atau seperti pelana. Sedangkan bab Hermitian berikutnya memainkan ulang seluruh
partiturnya atas $\C$.
\end{remark}

\section{Latihan}

\begin{exercise}[$\star$]\label{exo:b2:quadratic:1}
Reduksilah secara Gauss lalu berikan rank dan signaturnya:
\[
q_1(x,y) = x^2 + 4xy + y^2,
\qquad
q_2(x,y,z) = x^2 + 2y^2 + 3z^2 + 2xy + 2yz .
\]
\end{exercise}

\begin{exercise}[$\star$]\label{exo:b2:quadratic:2}
Diagonalkan secara ortogonal $A = \begin{pmatrix} 0 & 1 & 1\\ 1 & 0 &
1\\ 1 & 1 & 0\end{pmatrix}$ (dengan \omterm{def:b2:reduction:eigen}{nilai eigen} dari
perhitungan \cref{ch:b2:reduction}; kini buatlah basisnya
ortonormal) lalu reduksilah bentuk $q(x,y,z) = 2xy + 2yz + 2zx$ ke
sumbu utamanya.
\end{exercise}

\begin{exercise}[$\star$]\label{exo:b2:quadratic:3}
Buktikan bahwa $u^{**} = u$, $(u \circ v)^* = v^* \circ u^*$, dan bahwa
$\ker u^* = (\operatorname{im} u)^{\perp}$. Turunkan pula
$\operatorname{rk} u^* = \operatorname{rk} u$.
\end{exercise}

\begin{exercise}[$\star\star$]\label{exo:b2:quadratic:4}
Misalkan $A$ \omterm{def:b2:quadratic:adjoint}{simetrik} real dengan $A^3 = A$. Buktikan bahwa $A^2$
matriks sebuah proyeksi ortogonal. Lebih umum, kaitkanlah
dekomposisi spektral $A$ dan $P(A)$ untuk sebuah polinomial $P$.
\end{exercise}

\begin{exercise}[$\star\star$]\label{exo:b2:quadratic:5}
Buktikan bahwa $O(n) = \{P : P^{\mathsf T}P = I\}$ himpunan bagian yang
\omterm{def:b2:metric:compact}{kompak} pada $\mathcal{M}_n(\R)$ \emph{(tertutup: sebab prapeta $\{I\}$ oleh
pemetaan yang \omterm{def:b2:metric:continuity}{kontinu}; terbatas: sebab kolomnya vektor satuan)}. Apakah ia
\omterm{def:b2:metric:connected}{terhubung}?
\end{exercise}

\begin{exercise}[$\star\star$]\label{exo:b2:quadratic:6}
(Akar kuadrat) Misalkan $A$ \omterm{def:b2:quadratic:adjoint}{simetrik} semidefinit positif. Bangunlah
$B$ \omterm{def:b2:quadratic:adjoint}{simetrik} semidefinit positif dengan $B^2 = A$, lalu buktikan
bahwa ia \emph{tunggal} \emph{(keberadaannya: ambillah akar kuadrat
\omterm{def:b2:reduction:eigen}{nilai eigennya} dalam sebuah basis spektral; ketunggalannya: sebab calon $B$
komut dengan $A = B^2$, jadi memelihara \omterm{def:b2:reduction:eigen}{ruang eigennya} --- lalu reduksilah
ke kasus skalar pada masing-masingnya)}.
\end{exercise}

\begin{exercise}[$\star\star$]\label{exo:b2:quadratic:7}
Untuk $A$ \omterm{def:b2:quadratic:adjoint}{simetrik} real, buktikan bahwa $\vertiii{A}_2 :=
\sup_{\norm x_2 = 1}\norm{Ax}_2 = \max_i \abs{\lambda_i}$
(yakni jari-jari spektralnya), lalu hitunglah $\vertiii{A}_2$ untuk $A =
\begin{pmatrix} 1 & 2\\ 2 & 1\end{pmatrix}$.
\end{exercise}

\begin{exercise}[$\star\star\star$]\label{exo:b2:quadratic:8}
(Kriteria Sylvester) Misalkan $A$ \omterm{def:b2:quadratic:adjoint}{simetrik} real dengan minor
utama terdepan $\Delta_1, \dots, \Delta_n$ (yakni \omterm{def:b2:linalg:det}{determinan}
blok kiri atasnya). Buktikan bahwa $A$ definit positif bila dan hanya
bila semua $\Delta_k > 0$. \emph{(Untuk $\Rightarrow$: sebab pembatasan
bentuk definit bersifat definit, dan \omterm{def:b2:linalg:det}{determinan} matriks definit
positif --- yakni hasil kali \omterm{def:b2:reduction:eigen}{nilai eigennya} --- bersifat positif.
Untuk $\Leftarrow$: berinduksilah atas $n$; blok kiri atas berukuran $(n-1)$
bersifat definit positif, jadi diagonalkan bentuknya pada subruang itu lalu
lengkapkan kuadratnya dalam peubah terakhirnya; dan tanda entri diagonal
terakhirnya dikendalikan oleh $\det A = \Delta_n > 0$.)}
\end{exercise}

\begin{exercise}[$\star\star\star$]\label{exo:b2:quadratic:9}
(Courant--Fischer, \omterm{def:b2:reduction:eigen}{nilai eigen} kedua) Misalkan $u$ \omterm{def:b2:quadratic:adjoint}{simetrik} dengan
\omterm{def:b2:reduction:eigen}{nilai eigen} $\lambda_1 \geq \lambda_2 \geq \dots \geq \lambda_n$.
Buktikan
\[
\lambda_2 = \min_{\substack{H \text{ hiperbidang}}}\;
\max_{\substack{x \in H,\ \norm x = 1}} \langle u(x), x\rangle .
\]
\emph{(Untuk $\leq$: sebab sembarang hiperbidang memotong bidang-$2$ yang
direntang dua \omterm{def:b2:reduction:eigen}{vektor eigen} teratasnya. Untuk $\geq$: pilihlah $H =
(e_1)^{\perp}$.)}
\end{exercise}

\begin{exercise}[$\star\star$]\label{exo:b2:quadratic:10}
Tentukan rank dan signatur $q(x_1, \dots, x_n) = \sum_{i <
j} x_ix_j$ pada $\R^n$ (dengan $n \geq 2$), lewat dua cara: lewat kesamaan
aljabar $2q = \bigl(\sum x_i\bigr)^2 - \sum x_i^2$ bersama
pembatasan $q$ ke hiperbidang $\sum x_i = 0$;
dan lewat menghitung \omterm{def:b2:reduction:eigen}{nilai eigen} matriksnya $\frac12(J - I)$,
dengan $J$ menyatakan matriks serba-satu.
\end{exercise}

\begin{exercise}[$\star\star$]\label{exo:b2:quadratic:11}
Misalkan $A$ dan $B$ \omterm{def:b2:quadratic:adjoint}{simetrik} real dengan $B$ semidefinit positif.
Buktikan
\[
\lambda_{\min}(A)\operatorname{tr} B
\;\leq\;
\operatorname{tr}(AB)
\;\leq\;
\lambda_{\max}(A)\operatorname{tr} B .
\]
\emph{(Tulislah $B = C^{\mathsf T}C$ dan $\operatorname{tr}(AB) =
\sum_i \langle A c_i, c_i\rangle$ atas kolom $c_i$ milik
$C^{\mathsf T}$.)} Khususnya $\operatorname{tr}(AB) \geq 0$
bila keduanya semidefinit positif.
\end{exercise}

\begin{exercise}[$\star\star\star$]\label{exo:b2:quadratic:12}
Pada $E = \mathcal{M}_n(\R)$, tinjaulah $q(M) =
\operatorname{tr}(M^2)$.
\begin{enumerate}
  \item Tunjukkan bahwa $q$ \omterm{def:b2:quadratic:def}{bentuk kuadratik} dengan bentuk polar
        $\varphi(M, N) = \operatorname{tr}(MN)$.
  \item Tunjukkan bahwa matriks \omterm{def:b2:quadratic:adjoint}{simetrik} dan matriks antisimetrik
        membentuk subruang yang $\varphi$-ortogonal, dan di atasnya $q$
        berturut-turut definit positif dan definit negatif
        \emph{(hitunglah $\operatorname{tr}(M^2)$ entri demi entri pada
        masing-masing kasusnya)}.
  \item Simpulkan: $q$ bersignatur $\bigl(\frac{n(n+1)}{2},
        \frac{n(n-1)}{2}\bigr)$ dan berrank $n^2$.
\end{enumerate}
\end{exercise}

\section{Soal: Cholesky, Hadamard, dan dekomposisi
polar}

\begin{problem}\label{pb:b2:quadratic:1}
Teorema spektral itu mikroskop; sedangkan soal ini memakainya sebagai
pabrik. Dari matriks Gram kita memproduksi \emph{pemfaktoran
Cholesky}\index{pemfaktoran Cholesky} (sekaligus mengenali
pivot Gauss sebagai nisbah minor), lalu membuktikan \emph{ketaksamaan
Hadamard}\index{ketaksamaan Hadamard} tentang \omterm{def:b2:linalg:det}{determinan}, membangun
\emph{dekomposisi polar}\index{dekomposisi polar} $A =
QS$ dan \emph{dekomposisi nilai singular}, menggolongkan
konik bidang, lalu berakhir dengan reduksi serentak dua bentuk ---
yakni teorema di balik moda normal osilasi. Di sepanjang soal ini, $E =
\R^n$ dengan hasil kali dalam bakunya.

\medskip
\noindent\textbf{Bagian I --- Matriks Gram dan Cholesky.} Untuk
vektor $v_1, \dots, v_n \in E$, \emph{matriks Gram}-nya adalah
$G = \bigl(\langle v_i, v_j\rangle\bigr)_{i,j}$.
\begin{enumerate}
  \item Tunjukkan bahwa $G$ \omterm{def:b2:quadratic:adjoint}{simetrik} semidefinit positif, dan
        definit positif bila dan hanya bila $(v_1, \dots, v_n)$
        bebas linear \emph{(hitunglah $X^{\mathsf T}GX$)}.
  \item Sebaliknya, tunjukkan bahwa setiap $A$ \omterm{def:b2:quadratic:adjoint}{simetrik} semidefinit
        positif merupakan matriks Gram: $A = C^{\mathsf T}C$
        untuk suatu $C$ (pakailah akar kuadrat pada
        \cref{exo:b2:quadratic:6}), dengan $C$ terbalikkan bila dan hanya
        bila $A$ definit.
  \item Turunkan bahwa $A$ yang semidefinit positif memenuhi
        $\abs{a_{ij}} \leq \sqrt{a_{ii}\,a_{jj}}$ untuk setiap $i,
        j$ \emph{(batasilah ke dua koordinat)} --- yakni
        ketaksamaan Cauchy--Schwarz, dibaca ulang secara matriks.
  \item (Cholesky) Misalkan $A$ definit positif. Buktikan bahwa
        ada $T$ segitiga atas yang \emph{tunggal} dengan
        entri diagonal positif sedemikian sehingga
        \[
        A = T^{\mathsf T}\,T
        \]
        \emph{(keberadaannya: terapkan Gram--Schmidt pada vektor
        yang mewujudkan $A$ sebagai matriks Gram; ketunggalannya: bila
        $T_1^{\mathsf T}T_1 = T_2^{\mathsf T}T_2$, tunjukkan
        $T_1T_2^{-1}$ ortogonal sekaligus segitiga dengan
        diagonal positif, sehingga $I$)}.
  \item Tunjukkan bahwa minor utama terdepannya memenuhi
        $\Delta_k = (t_{11}\cdots t_{kk})^2$, lalu turunkan bahwa
        pivot yang dihasilkan \omterm{thm:b2:quadratic:gauss}{reduksi Gauss} sebuah
        bentuk definit positif, diambil menurut urutan peubah
        alaminya, adalah
        \[
        d_k = \frac{\Delta_k}{\Delta_{k-1}}
        \qquad (\Delta_0 = 1) :
        \]
        jadi minor pada kriteria Sylvester
        (\cref{exo:b2:quadratic:8}) dan pivot Gauss adalah
        data yang sama. Periksalah pada
        \cref{ex:b2:quadratic:tworoads}.
\end{enumerate}

\medskip
\noindent\textbf{Bagian II --- \omterm{pb:b2:quadratic:1}{Ketaksamaan Hadamard}.}
\begin{enumerate}[resume]
  \item Misalkan $A$ definit positif. Buktikan
        \[
        \det A \leq a_{11}\,a_{22}\cdots a_{nn}
        \]
        \emph{(normalkan: $B = DAD$ dengan $D =
        \operatorname{diag}(a_{ii}^{-1/2})$ berdiagonal satuan;
        lalu batasi $\det B = \prod \mu_i$ lewat AM--GM terhadap
        $\operatorname{tr} B = n$)}.
  \item Tunjukkan bahwa kesamaannya berlaku bila dan hanya bila $A$ diagonal.
  \item Turunkan \emph{\omterm{pb:b2:quadratic:1}{ketaksamaan Hadamard}}: untuk setiap matriks
        real bujur sangkar $M$ berkolom $c_1, \dots, c_n$,
        \[
        \abs{\det M} \leq \prod_{i=1}^{n}\norm{c_i}_2 ,
        \]
        dengan kesamaannya (bagi $M$ yang terbalikkan) berlaku bila dan
        hanya bila kolomnya ortogonal berpasangan \emph{(terapkan
        pertanyaan 6--7 pada $M^{\mathsf T}M$)}.
  \item Panen geometri dan kombinatorikanya: tafsirkanlah
        pertanyaan 8 sebagai ``volume sebuah paralelepipedum paling
        banyak sebesar hasil kali panjang rusuknya''; lalu tunjukkan bahwa
        matriks yang semua entrinya di $\intcc{-1}{1}$ punya
        $\abs{\det M} \leq n^{n/2}$. (Matriks yang mencapai batas
        ini --- yakni matriks Hadamard --- ada untuk $n = 1, 2$
        dan banyak kelipatan $4$; sedangkan apakah ada untuk \emph{setiap}
        kelipatan $4$ merupakan masalah \omterm{def:b2:metric:topology}{terbuka} yang terkenal.)
\end{enumerate}

\medskip
\noindent\textbf{Bagian III --- \omterm{pb:b2:quadratic:1}{Dekomposisi polar} dan nilai
singular.}
\begin{enumerate}[resume]
  \item Misalkan $A$ terbalikkan. Tunjukkan bahwa $A^{\mathsf T}A$
        bersifat definit positif, dan bahwa
        \[
        S = \sqrt{A^{\mathsf T}A}
        \quad\text{(yakni akar kuadrat pada
        \cref{exo:b2:quadratic:6})},
        \qquad Q = AS^{-1}
        \]
        memberi pemfaktoran $A = QS$ dengan $Q$ ortogonal dan
        $S$ definit positif.
  \item Buktikan bahwa pemfaktoran $A$ yang terbalikkan itu bersifat
        tunggal.
  \item Perluaslah keberadaannya ke sembarang $A$: pilihlah
        $\varepsilon_k \to 0$ dengan $A + \varepsilon_k I$
        terbalikkan, tulislah $A + \varepsilon_kI = Q_kS_k$, lalu
        pakailah kekompakan $O(n)$
        (\cref{exo:b2:quadratic:5}) untuk menyarikan $Q_k \to Q$;
        lalu tunjukkan $S_k = Q_k^{\mathsf T}(A + \varepsilon_kI)$
        konvergen ke suatu $S$ yang semidefinit positif dengan $A =
        QS$ dan $S = \sqrt{A^{\mathsf T}A}$. Di mana
        ketunggalannya gagal untuk $A$ yang singular?
  \item (Dekomposisi nilai singular) Turunkan bahwa setiap $A$
        real bujur sangkar dituliskan
        \[
        A = U\,\Sigma\,V^{\mathsf T},
        \qquad U, V \in O(n),\quad
        \Sigma = \operatorname{diag}(\sigma_1, \dots,
        \sigma_n),\ \sigma_i \geq 0 ,
        \]
        dengan $\sigma_i$ (yakni \emph{nilai singularnya}) adalah
        \omterm{def:b2:reduction:eigen}{nilai eigen} $\sqrt{A^{\mathsf T}A}$.
  \item Tiga akibatnya: $\vertiii{A}_2 = \sigma_{\max}$ bagi
        \emph{setiap} $A$ real (yang menyamaratakan
        \cref{exo:b2:quadratic:7}); lalu $\abs{\det A} =
        \sigma_1\cdots\sigma_n$; dan peta bola satuan
        oleh $A$ yang terbalikkan adalah sebuah elipsoid dengan
        setengah sumbu $\sigma_1, \dots, \sigma_n$ sepanjang kolom
        $U$.
\end{enumerate}

\medskip
\noindent\textbf{Bagian IV --- Konik, lewat teorema spektral.}
Sebuah konik bidang adalah himpunan nol $f(x) = q(x) + \langle b,
x\rangle + c$, dengan $q \neq 0$ \omterm{def:b2:quadratic:def}{bentuk kuadratik} bermatriks $A$,
lalu $b \in \R^2$ dan $c \in \R$.
\begin{enumerate}[resume]
  \item Reduksilah $f$ lewat sebuah perputaran (yakni sumbu utamanya,
        \cref{cor:b2:quadratic:principalaxes}) lalu sebuah
        penggeseran, kemudian golongkanlah bentuk tak kosong dan
        tak merosot yang mungkin, menurut signatur $q$: elips
        (bila $\det A > 0$), hiperbola (bila $\det A < 0$), parabola
        (bila $\det A = 0$, berrank $1$, dengan suku linearnya tak
        terserap).
  \item Jalankanlah reduksinya selengkapnya untuk
        \[
        x^2 + 4xy + y^2 + 2x - 2y = 4 :
        \]
        yakni koordinat terputarnya, persamaan tereduksinya, lalu sifat dan
        pusat koniknya.
  \item (Konik berpusat) Andaikan $\det A \neq 0$. Tunjukkan bahwa
        \emph{pusatnya} adalah $x_0 = -\frac12 A^{-1}b$, dan bahwa
        \omterm{def:b2:quadratic:def}{kekongruenan} oleh $\begin{pmatrix} I & x_0\\ 0 &
        1\end{pmatrix}$ atas matriks $3\times3$ $\widetilde Q
        = \begin{pmatrix} A & b/2 \\ b^{\mathsf T}/2 &
        c\end{pmatrix}$ menghasilkan
        \[
        \det\widetilde Q = f(x_0)\,\det A :
        \]
        konik berpusatnya merosot (menjadi sebuah titik atau dua garis)
        persis ketika $\det\widetilde Q = 0$.
  \item Periksalah pertanyaan 17 pada contoh pertanyaan 16:
        hitunglah $x_0$, $f(x_0)$ dan $\det\widetilde Q$, lalu
        simpulkan sekali lagi bahwa koniknya adalah
        hiperbola yang tak merosot.
  \item (Sebuah berkas kuadrik) Untuk $\lambda \in \R$, golongkanlah
        permukaan
        \[
        x^2 + y^2 + z^2 + 2\lambda(xy + yz + zx) = 1
        \]
        lewat \omterm{def:b2:reduction:eigen}{nilai eigen} matriksnya \emph{(dengan struktur
        serba-satu: \omterm{def:b2:reduction:eigen}{nilai eigennya} $1 + 2\lambda$ dan $1 -
        \lambda$ yang ganda)}: bola/elipsoid, silinder, sepasang
        bidang, hiperboloid berdaun satu dan berdaun dua, menurut
        nilai $\lambda$.
\end{enumerate}

\medskip
\noindent\textbf{Bagian V --- Dua bentuk sekaligus: reduksi
serentak.}
\begin{enumerate}[resume]
  \item Misalkan $q$ definit positif dan $q'$ sembarang
        \omterm{def:b2:quadratic:def}{bentuk kuadratik} pada $E$. Buktikan bahwa ada sebuah basis
        $E$ yang ortonormal terhadap $q$ dan ortogonal terhadap
        $q'$: di dalamnya $q = \sum x_i^2$ dan $q' = \sum \mu_i
        x_i^2$ \emph{(pakailah $q$ sebagai hasil kali dalamnya lalu terapkan
        teorema spektral pada endomorfisma yang mewakili
        $q'$)}.
  \item Bentuk matriksnya: untuk $A$ definit positif dan $B$
        \omterm{def:b2:quadratic:adjoint}{simetrik}, ada $P$ terbalikkan dengan
        $P^{\mathsf T}AP = I$ dan $P^{\mathsf T}BP =
        \operatorname{diag}(\mu_1, \dots, \mu_n)$, dengan
        $\mu_i$ menyatakan akar $\det(B - \mu A) = 0$.
  \item Jalankanlah selengkapnya untuk
        \[
        A = \begin{pmatrix} 2 & 1\\ 1 & 1 \end{pmatrix},
        \qquad
        B = \begin{pmatrix} 0 & 1\\ 1 & 0 \end{pmatrix} :
        \]
        yakni \omterm{def:b2:reduction:eigen}{nilai eigen} tergeneralisasinya $\mu_\pm$, beserta vektor
        yang mendiagonalkan kedua bentuknya.
  \item Tunjukkan bahwa kedefinitan positifnya tak dapat dilepas: untuk
        \[
        A = \begin{pmatrix} 1 & 0\\ 0 & -1 \end{pmatrix},
        \qquad
        B = \begin{pmatrix} 0 & 1\\ 1 & 0 \end{pmatrix},
        \]
        tak ada basis yang mendiagonalkan kedua bentuknya \emph{(sebab bila $P$
        mendiagonalkan keduanya, $\det(B - \mu A)$ akan terurai dengan
        akar yang real; hitunglah ia)}.
  \item Tunjukkan bahwa $\mu_i$ pada pertanyaan 21 adalah
        \omterm{def:b2:reduction:eigen}{nilai eigen} $A^{-1}B$, dan bahwa $A^{-1}B$, walau tak
        \omterm{def:b2:quadratic:adjoint}{simetrik} secara umum, selalu \omterm{def:b2:reduction:diag}{dapat didiagonalkan} dengan
        \omterm{def:b2:reduction:eigen}{nilai eigen} yang real \emph{(konjugasikan dengan $\sqrt A$)}.
  \item Rangkuman. Satu kalimat untuk masing-masing: (i) satu teorema
        yang disandari setiap Bagiannya; (ii) hasil Bagian
        I--III mana yang bertahan untuk matriks \emph{semi}definit
        positif, dan mana yang menuntut kedefinitan; (iii) sistem
        fisis yang osilasi kecilnya didiagonalkan pertanyaan
        20--22 (dengan energi kinetik dan energi potensial sebagai
        kedua bentuknya), dan apa makna $\mu_i$ di sana.

\end{enumerate}
\end{problem}
